\documentclass[12pt]{article}
\usepackage[utf8]{inputenc}
\usepackage{hyperref}
\usepackage{amsmath}
\usepackage{geometry}
\geometry{margin=1in}

\title{Statistics Project: Central Limit Theorem Simulation}
\author{}
\date{}

\begin{document}
\maketitle

\section*{Interactive Simulation}

This report is accompanied by an interactive browser-based simulation that
demonstrates the \textbf{Central Limit Theorem} (CLT).  The simulation lets
you choose a population distribution (Uniform, Exponential, Skewed, or
Bimodal), set the sample size~$n$ and the number of repeated samples, and
observe the distribution of sample means converge to a normal distribution
as $n$ grows.

\medskip
\noindent
To launch the simulation, click the link below (requires an internet
connection):

\medskip
\noindent
\href{https://subzeroxx5.github.io/Project-for-Statistics/simulation.html}{Click to launch}

\section*{Background}

The Central Limit Theorem states that, given a population with mean
$\mu$ and finite variance $\sigma^{2}$, the sampling distribution of
the sample mean $\bar{X}_{n}$ approaches a normal distribution as the
sample size $n \to \infty$:

\[
  \bar{X}_{n} \;\xrightarrow{d}\; \mathcal{N}\!\left(\mu,\,\frac{\sigma^{2}}{n}\right).
\]

The standard error of the mean is therefore $\sigma / \sqrt{n}$, which
the simulation reports alongside the empirically observed standard
deviation of the sample means.

\end{document}
